India astronomy satellite, AstroSat, launched in September 2015, has five
different instruments.

\ioaafig{2016_TH_T12-f1.pdf}

In this question, we will discuss three of these instruments (SXT, LAXPC,
CZTI), which point in the same direction and observe in X-ray wavelengths. The
details of these instruments are given in the table below.

\begin{center}
\small
\begin{tabular}{|l|c|c|c|p{3.2cm}|c|}
\hline
Instrument & Band & Collecting & Effective Photon & Saturation level & No.\ of
\\
 & [keV] & Area [m$^2$] & Detection Efficiency & [counts] & Pixels \\
\hline
SXT & 0.3--80 & 0.067 & 60\% & 15000 (total) & $512 \times 512$ \\
\hline
LAXPC & 3--80 & 1.5 & 40\% & 50000 (in any one counter) or 200000 (total)
  & --- \\
\hline
CZTI & 10--150 & 0.09 & 50\% & --- & $4 \times 4096$ \\
\hline
\end{tabular}
\end{center}

You should note that LAXPC energy range is divided into 8 different energy band
counters of equal bandwidth with no overlap.

\begin{description}
\item[(T12.1)] Some X-ray sources like Cas A have a prominent emission line at
  0.01825\,nm corresponding to a radioactive transition of $^{44}$Ti. Suppose
  there exists a source which emits only one bright emission line corresponding
  to this transition. What should be the minimum relative velocity ($v$) of the
  source, which will make the observed peak of this line to get registered in a
  different energy band counter of LAXPC as compared to a source at rest?
\end{description}

These instruments were used to observe an X-ray source (assumed to be a point
source), whose energy spectrum followed the power law,
\[ F(E) = K E^{-2/3} \qquad
   [\text{in units of counts/keV/m}^2\text{/s}] \]
where $E$ is the energy in keV, $K$ is a constant and $F(E)$ is photon flux
density at that energy. Photon flux density, by definition, is given for per
unit collecting area (m$^2$) per unit bandwidth (keV) and per unit time
(seconds). From prior observations, we know that the source has a flux density
of 10\,counts/keV/m$^2$/s at 1\,keV, when measured by a detector with 100\%
photon detection efficiency. The ``counts'' here mean the number of photons
reported by the detector.

As the source flux follows the power law given above, we know that for a given
energy range from $E_1$ (lower energy) to $E_2$ (higher energy) the total
photon flux ($F_{\mathrm{T}}$) will be given by
\[ F_{\mathrm{T}} = 3K\left(E_2^{1/3} - E_1^{1/3}\right) \qquad
   [\text{in units of counts/m}^2\text{/s}] \]

\begin{description}
\item[(T12.2)] Estimate the incident flux density from the source at 1\,keV,
  5\,keV, 40\,keV and 100\,keV. Also estimate what will be the total count per
  unit bandwidth recorded by each of the instruments at these energies for an
  exposure time of 200 seconds.
\item[(T12.3)] For this source, calculate the maximum exposure time
  ($t_{\mathrm{S}}$), without suffering from saturation, for the CCD of SXT.
\item[(T12.4)] If the source became 3500 times brighter, calculate the expected
  counts per second in LAXPC counter 1, counter 8 as well as total counts
  across the entire energy range. If we observe for longer period, will the
  counter saturate due to any individual counter or due to the total count?
  Tick the appropriate box in the Summary Answersheet.
\item[(T12.5)] Assume that the counts reported by CZTI due to random
  fluctuations in electronics are about 0.00014 counts per pixel per keV per
  second at all energy levels. Any source is considered as ``detected'' when
  the SNR (signal to noise ratio) is at least 3. What is minimum exposure time,
  $t_c$, needed for the source above to be detected in CZTI?\\
  Note that the ``noise'' in a detector is equal to the square root of the
  counts due to random fluctuations.
\item[(T12.6)] Let us consider the situation where the source shows variability
  in number flux, so that the factor $K$ increases by 20\%. AstroSat observed
  this source for 1 second before the change and 1 second after this change in
  brightness. Calculate the counts measured by SXT, LAXPC and CZTI in both the
  observations. Which instrument is best suited to detect this change? Tick the
  appropriate box in the Summary Answersheet.
\end{description}
